\question[7]
L'\textbf{ASLR} (Address Space Layout Randomization). (a) Che cosa randomizza e perché ostacola l'attaccante? (b) Quali sono i tre livelli su Linux e perché serve anche \textbf{PIE}? (c) Elenca le principali tecniche per aggirarlo, indicando quella decisiva.
\begin{solution}
(a) Posiziona in modo casuale, a ogni esecuzione: l'indirizzo base dell'eseguibile, le librerie condivise, l'heap e lo stack. Se l'attaccante non conosce gli indirizzi, non può puntare in modo affidabile allo shellcode o alle funzioni da riusare: trasforma indirizzi ``fissi e noti'' in ``imprevedibili''.\\
(b) \lstinline|/proc/sys/kernel/randomize_va_space|: \textbf{0} nessuna randomizzazione; \textbf{1} conservativa (stack, mmap, librerie, \emph{non} l'eseguibile); \textbf{2} completa (anche l'heap e, con PIE, l'eseguibile). Per singolo processo: \lstinline|setarch -R| / \lstinline|personality(ADDR_NO_RANDOMIZE)|. L'\textbf{eseguibile} è randomizzabile \emph{solo} se compilato \textbf{PIE}; con \lstinline|-no-pie| il codice del binario resta a indirizzo \emph{fisso} (tipicamente \lstinline|0x0804...| a 32 bit) --- come i binari del laboratorio.\\
(c) \textbf{Information leak} (decisivo): far rivelare al programma un indirizzo reale a runtime; un solo indirizzo di libc leakato $\to$ si ricavano tutti gli altri (offset interni fissi) e si costruisce al volo ROP/ret2libc. Altre: \textbf{forza bruta + NOP sled} (a 32 bit l'entropia è bassa; il NOP sled allarga la ``zona di atterraggio''; a 64 bit impraticabile); \textbf{partial overwrite} (sovrascrivere solo i byte bassi, che l'ASLR non cambia, per saltare ``vicino''); side-channel / de-randomizzazione. Morale: l'ASLR alza il costo, da solo non impedisce l'attacco.
\end{solution}
